% 2009 年辽宁卷文科第 10 题：计算月总收入 S 与月净盈利 V 的程序框图。
% 来源 PDF：普通高考/2009/2009辽宁文.pdf，SHA256 99623efa457632f26bc5d5775293bf9fd13a56add984179360e30698afc7cb73
% 原页：第 1 页（A3 三栏）第 2 栏上部，印刷页 539。
% 重绘依据：开始 → 输入 N,a1,...,aN → k=1,S=0,T=0 → A=ak → 判断框（原页空白）：
%   否（左）T=T+A，是（右）S=S+A；两路汇合后判断 k<N：是（右）k=k+1 回到 A=ak，
%   否（下）输出 S,V → 结束。判断框内的条件在原页留空，由选项填入。
\begin{tikzpicture}[scale=0.82, semithick, >=Stealth,
  every node/.style={inner sep=3pt, font=\small}]
  \node (start) at (0,0) [draw, rounded corners=4pt, minimum width=2.1cm, minimum height=0.6cm] {开始};
  \node (input) at (0,-1.3) [draw, trapezium, trapezium left angle=75, trapezium right angle=105,
    minimum width=4.9cm, minimum height=0.72cm] {输入 $N,\,a_1,\,a_2,\,\cdots,\,a_N$};
  \node (init) at (0,-2.75) [draw, minimum width=4.0cm, minimum height=0.7cm] {$k=1,\ S=0,\ T=0$};
  \node (ak) at (0,-4.05) [draw, minimum width=1.75cm, minimum height=0.62cm] {$A=a_k$};
  \node (cond) at (0,-5.55) [draw, diamond, aspect=2.4, minimum width=3.3cm, minimum height=1.1cm] {};
  \node (t) at (-2.75,-7.1) [draw, minimum width=2.0cm, minimum height=0.62cm] {$T=T+A$};
  \node (s) at (2.75,-7.1) [draw, minimum width=2.0cm, minimum height=0.62cm] {$S=S+A$};
  \node (kn) at (0,-8.9) [draw, diamond, aspect=2.2, minimum width=2.5cm, minimum height=1.0cm] {$k<N$};
  \node (inc) at (5.0,-8.05) [draw, minimum width=2.0cm, minimum height=0.62cm] {$k=k+1$};
  \node (out) at (0,-10.55) [draw, trapezium, trapezium left angle=75, trapezium right angle=105,
    minimum width=2.7cm, minimum height=0.68cm] {输出 $S,\,V$};
  \node (end) at (0,-11.8) [draw, rounded corners=4pt, minimum width=2.0cm, minimum height=0.6cm] {结束};
  \coordinate (merge) at (0,-8.0);

  \draw[->] (start) -- (input);
  \draw[->] (input) -- (init);
  \draw[->] (init) -- (ak);
  \draw[->] (ak) -- (cond);
  \draw[->] (cond.west) -| node[pos=0.28, above] {否} (t.north);
  \draw[->] (cond.east) -| node[pos=0.28, above] {是} (s.north);
  \draw (t.south) -- (t.south |- merge) -- (merge);
  \draw (s.south) -- (s.south |- merge) -- (merge);
  \draw[->] (merge) -- (kn.north);
  \draw[->] (kn.east) -- node[below=2pt, pos=0.45] {是} (inc.west);
  \draw[->] (inc.north) |- (ak.east);
  \draw[->] (kn.south) -- node[right=2pt, pos=0.45] {否} (out.north);
  \draw[->] (out) -- (end);
\end{tikzpicture}
